\documentclass[11pt,a4paper]{article}
\usepackage[utf8]{inputenc}
\usepackage[T1]{fontenc}
\usepackage[brazil]{babel}
\usepackage{lmodern}
\usepackage[margin=2cm]{geometry}
\usepackage{amsmath,amssymb}
\usepackage{booktabs}
\usepackage{graphicx}
\usepackage{hyperref}
\usepackage{xcolor}
\usepackage{titlesec}
\usepackage{tabularx}
\usepackage{enumitem}
\usepackage{float}
\usepackage{array}
\renewcommand{\tabularxcolumn}[1]{>{\raggedright\arraybackslash}p{#1}}
\usepackage{newunicodechar}
\usepackage[most]{tcolorbox}
\usepackage{etoolbox}

\hypersetup{colorlinks=true, linkcolor=blue!70!black, citecolor=blue!70!black, urlcolor=blue!70!black}

\titleformat{\section}{\large\bfseries}{\thesection.}{0.5em}{}
\titleformat{\subsection}{\normalsize\bfseries}{\thesubsection.}{0.5em}{}
\titlespacing*{\section}{0pt}{1.2ex plus 0.5ex minus .2ex}{0.6ex}
\titlespacing*{\subsection}{0pt}{1.0ex plus 0.4ex minus .2ex}{0.4ex}
\setlength{\parskip}{0.45em}
\setlength{\parindent}{0pt}
\setlist{itemsep=0.15em,topsep=0.3em}

% --- símbolos ---
\newunicodechar{π}{\ensuremath{\pi}}
\newunicodechar{τ}{\ensuremath{\tau}}
\newunicodechar{φ}{\ensuremath{\varphi}}
\newunicodechar{ψ}{\ensuremath{\psi}}
\newunicodechar{Σ}{\ensuremath{\Sigma}}
\newunicodechar{→}{\ensuremath{\rightarrow}}
\newunicodechar{⇒}{\ensuremath{\Rightarrow}}
\newunicodechar{∈}{\ensuremath{\in}}
\newunicodechar{−}{\ensuremath{-}}
\newunicodechar{≈}{\ensuremath{\approx}}
\newunicodechar{≠}{\ensuremath{\neq}}
\newunicodechar{≡}{\ensuremath{\equiv}}
\newunicodechar{≤}{\ensuremath{\leq}}
\newunicodechar{≥}{\ensuremath{\geq}}
\newunicodechar{′}{\ensuremath{'}}
\newunicodechar{‑}{-}
\newunicodechar{✔}{\ensuremath{\checkmark}}
\newunicodechar{✘}{\ensuremath{\times}}

% --- cores e caixas ---
\definecolor{cC}{HTML}{7A4FA3}
\definecolor{cS}{HTML}{2F7D74}
\definecolor{cAcc}{HTML}{2F5D8A}
\newcommand{\metrictag}[2]{{\setlength{\fboxsep}{1.5pt}%
  \ifstrequal{#1}{c}{\colorbox{cC}{\textcolor{white}{\scriptsize\textbf{#2}}}}{%
  \ifstrequal{#1}{s}{\colorbox{cS}{\textcolor{white}{\scriptsize\textbf{#2}}}}{%
  \colorbox{black!60}{\textcolor{white}{\scriptsize\textbf{#2}}}}}}}
\newcommand{\pagechip}[1]{{\setlength{\fboxsep}{1.5pt}\colorbox{cAcc}{\textcolor{white}{\scriptsize\texttt{#1}}}}}
\newcommand{\readtag}[1]{\texttt{\small #1}}
\newcommand{\partbar}[1]{\par\vspace{1.2em}\noindent\colorbox{black!80}{\parbox{\dimexpr\linewidth-2\fboxsep}{\color{white}\bfseries\sffamily\MakeUppercase{#1}}}\par\vspace{0.4em}}
\newtcolorbox{citacard}[1]{breakable, enhanced, colback=white, colframe=black!25, boxrule=0.5pt, arc=2pt,
  left=6pt, right=6pt, top=4pt, bottom=4pt, title={\small #1}, coltitle=black, colbacktitle=black!6, fonttitle=\normalfont}
\newenvironment{origquote}{\par\itshape\small}{\par}
\newenvironment{tradquote}{\par\small\textsc{Tradução:}\ }{\par}
\newcommand{\explab}[1]{\par\smallskip{\small\color{cAcc}\textbf{\MakeUppercase{#1}}}\par}
\colorlet{notebg}{blue!4}\colorlet{notefg}{cAcc}
\colorlet{okbg}{green!5}\colorlet{okfg}{green!45!black}
\colorlet{warnbg}{orange!8}\colorlet{warnfg}{orange!80!black}
\newtcolorbox{infobox}[2][note]{breakable, enhanced, boxrule=0pt, arc=2pt, left=6pt, right=6pt, top=4pt, bottom=4pt,
  colback=#1bg, borderline west={3pt}{0pt}{#1fg}, frame hidden,
  before upper={\ifstrempty{#2}{}{\textbf{#2}\par}}}
\colorlet{pairc}{cC}\colorlet{pairs}{cS}
\newenvironment{paircol}[2]{\begin{tcolorbox}[width=0.485\linewidth, nobeforeafter, box align=top, boxrule=0.5pt,
  colframe=black!25, colback=white, borderline north={3pt}{0pt}{pair#1}, left=5pt, right=5pt,
  before upper={{\color{pair#1}\small\textbf{\MakeUppercase{#2}}}\par}]}{\end{tcolorbox}}

\newunicodechar{ℓ}{\ensuremath{\ell}}
\newunicodechar{⊙}{\ensuremath{\odot}}
\newunicodechar{α}{\ensuremath{\alpha}}
\newunicodechar{β}{\ensuremath{\beta}}

\newunicodechar{×}{\ensuremath{\times}}
\newunicodechar{±}{\ensuremath{\pm}}

\begin{document}

\begin{center}
    {\LARGE \textbf{Fichamento: XAI para auditoria de modelos de ECG}} \\[0.35em]
    {\large WAGNER, P.; MEHARI, T.; HAVERKAMP, W.; STRODTHOFF, N. \textit{Explaining deep learning for ECG analysis: Building blocks for auditing and knowledge discovery}. Computers in Biology and Medicine, v.~176, 108525, 2024 (arXiv:2305.17043v2)} \\[0.3em]
    \textbf{Tema:} testes de sanidade por regressão de amplitudes, atribuições agregadas por batimento e segmento (“glocal”), TCAV e descoberta de subgrupos \\[0.2em]
    \textbf{Aluno:} Arthur Santos de Oliveira (RA: 156297) \\[0.2em]
    \textbf{Atualizado em:} 28/10/2026 \\[0.2em]
    \textbf{Link:} \url{https://doi.org/10.1016/j.compbiomed.2024.108525} \quad (\url{https://arxiv.org/abs/2305.17043}) \\[0.2em]
    \textbf{Leitura:} íntegra do texto principal (arXiv v2, 25 páginas); apêndices (segmentação, TCAV técnico, figuras suplementares) apenas percorridos \\[0.2em]
    \textbf{Repositório:} \url{https://github.com/Arthur-so/shap-ecg-arrhythmia}
\end{center}
\vspace{-0.5em}
\hrule
\vspace{0.8em}

\textbf{Como ler.} \texttt{[p. N]} indica a página do PDF lido (arXiv v2, 25 páginas, incluindo apêndices). As citações estão no original em inglês, cada uma seguida de uma tradução livre e de uma explicação.

\begin{infobox}[note]{Resumo}
\begin{itemize}
  \item \textbf{Objetivo:} propor “blocos” de XAI para auditar modelos de ECG antes do uso clínico e para descobrir conhecimento, indo além de exemplos escolhidos a dedo.
  \item \textbf{Dados e modelos:} PTB-XL (23 rótulos subdiagnósticos; 100 Hz; janelas de 2,5 s; multirrótulo com sigmoide e entropia cruzada binária); LeNet rasa e xresnet1d50.
  \item \textbf{Métodos de explicação:} Grad-CAM, \textit{saliency} (norma do gradiente), Integrated Gradients (IG) e LRP; depois, só \textit{saliency}; TCAV (vetores de ativação de conceitos) na análise global.
  \item \textbf{Avaliação:} teste de sanidade com modelos de regressão das amplitudes de P, R e T em cada derivação: mede-se a especificidade espacial (fração da atribuição na derivação-alvo) e temporal (perfil médio sobre os batimentos).
  \item \textbf{Conclusão:} só a \textit{saliency} passa; Grad-CAM, IG e LRP concentram relevância no QRS mesmo quando a tarefa é a onda P ou T. Mapas agregados por segmento concordam com critérios clínicos (LVH, BRE completo, infarto).
\end{itemize}
\end{infobox}

\section{O que o trabalho faz}
O artigo tem quatro partes. (1) \textbf{Testes de sanidade}: treinam-se modelos que regridem, a partir do ECG, a amplitude da onda P, do pico R ou da onda T em cada derivação (12 saídas, valores do programa de Glasgow no PTB-XL+). Um método sensato deveria atribuir à derivação e ao segmento correspondentes. (2) \textbf{XAI “glocal”}: as atribuições locais são alinhadas por batimento (300 ms antes e 500 ms depois do pico R) ou ponderadas por uma segmentação probabilística em 24 segmentos (U-Net 2D própria) e agregadas sobre as 100 maiores predições de cada classe, dando mapas por derivação $\times$ segmento. (3) \textbf{XAI global}: TCAV testa se conceitos clínicos (QRS largo para BRE, índice de Sokolow-Lyon para LVH, ondas Q para infarto, idade, sexo) são usados pelo modelo, com 10 CAVs $\times$ 10 modelos. (4) \textbf{Descoberta}: agrupam-se atribuições alinhadas para separar subclasses de infarto.

\section{Como avalia as explicações}
A avaliação quantitativa dos métodos está só no teste de sanidade. A especificidade espacial $s_l$ é a razão entre a norma temporal da atribuição na derivação-alvo e a soma das normas nas 12 derivações (1 = toda a atribuição na derivação certa). A especificidade temporal recorta cada batimento (80 amostras), tira a norma entre derivações e faz a média entre batimentos; espera-se um pico no segmento regredido. Não há perturbação, valor de referência nem métrica de fidelidade do tipo deleção: o critério é uma “verdade de referência” construída pela própria tarefa de regressão. Na parte glocal, a concordância com critérios clínicos é qualitativa (os sete segmentos de maior relevância por classe).

\section{Resultados e conclusões}
\begin{itemize}
  \item \textbf{Sanidade:} LRP e IG são os mais específicos nas precordiais; a \textit{saliency}, nas derivações dos membros, com menor variância. No eixo temporal, IG, LRP e Grad-CAM põem a maior relevância no pico R mesmo para P e T; só a \textit{saliency} tem o maior valor no segmento correto nos seis experimentos.
  \item \textbf{Glocal:} LVH concentra relevância nos picos R de V5, V1 e V6; BRE completo em V1; infarto anterior em V1--V2 (QRS) e inferior em II, III, aVF e aVL (ondas Q).
  \item \textbf{Global:} os três conceitos clínicos são usados de forma consistente pelos dois modelos para as classes correspondentes.
  \item \textbf{Descoberta:} separar AMI de IMI agrupando atribuições dá acurácia de 90\%, contra 75\% com os batimentos de entrada e 61\% com as ativações internas.
  \item \textbf{Limitações declaradas:} só o PTB-XL; uma única atribuição (\textit{saliency}) nas análises seguintes; a análise glocal só capta padrões iguais em todos os batimentos (não capta variação de RR).
\end{itemize}

\section{Citações}

\begin{citacard}{\pagechip{p. 8} §3.1 Testes de sanidade}
\begin{origquote}“In terms of temporal specificity, IG and LRP attribute the largest fraction of relevance to the R-peak even if the R-peak is not directly relevant for the given task (e.g., for predicting the P-wave or T-wave amplitude). Interestingly, a similar effect is observed for GradCAM.”\end{origquote}
\begin{tradquote}“Em termos de especificidade temporal, IG e LRP atribuem a maior fração da relevância ao pico R mesmo quando ele não é diretamente relevante para a tarefa (por exemplo, prever a amplitude da onda P ou T). Curiosamente, efeito semelhante é observado no Grad-CAM.”\end{tradquote}
\explab{Explicação}
É evidência direta, em ECG, do viés de amplitude: métodos do tipo gradiente $\times$ entrada concentram relevância no QRS. O GradientSHAP do TCC pertence a essa família (multiplica o gradiente por ECG $-$ referência), e o achado do eMPRT (o mapa do modelo de pesos sorteados também se concentra no QRS) é coerente com isso.
\end{citacard}

\begin{citacard}{\pagechip{p. 14} §4 Discussão}
\begin{origquote}“Surprisingly, three out of the four considered attribution methods do not show a sensible degree of temporal or spatial specificity and show a strong bias toward attributing relevance to the QRS-complex, which typically contains the highest signal amplitudes.”\end{origquote}
\begin{tradquote}“Surpreendentemente, três dos quatro métodos de atribuição considerados não mostram grau razoável de especificidade temporal ou espacial e têm forte viés de atribuir relevância ao complexo QRS, que costuma conter as maiores amplitudes do sinal.”\end{tradquote}
\explab{Explicação}
Dá nome ao risco que o TCC precisa controlar: um mapa que destaca o QRS pode refletir a amplitude, e não a decisão. Por isso o TCC compara as distâncias com um piso e com um modelo aleatório, em vez de ler a concentração no QRS como sinal de qualidade.
\end{citacard}

\begin{citacard}{\pagechip{p. 14} §4 Discussão}
\begin{origquote}“Gradient noise, a major drawback of saliency methods, can be very effectively reduced by aggregating attributions across beats and samples.”\end{origquote}
\begin{tradquote}“O ruído do gradiente, grande desvantagem dos métodos de \textit{saliency}, pode ser reduzido de forma muito eficaz agregando as atribuições entre batimentos e amostras.”\end{tradquote}
\explab{Explicação}
Apoia a agregação usada no TCC: somar a atribuição dentro de regiões (pré-QRS, QRS, ST-T por derivação) reduz a variação ponto a ponto, que no GradientSHAP é inflada pelo sorteio de referências e ruído (o “piso”).
\end{citacard}

\begin{citacard}{\pagechip{p. 8} §3.2 XAI glocal}
\begin{origquote}“It is worth stressing, that while local attributions can potentially aid human decision-making, they are not directly applicable for auditing, which focuses on model behavior across entire patient groups.”\end{origquote}
\begin{tradquote}“Vale ressaltar que, embora as atribuições locais possam ajudar a decisão humana, elas não se aplicam diretamente à auditoria, que se ocupa do comportamento do modelo em grupos inteiros de pacientes.”\end{tradquote}
\explab{Explicação}
Justifica resumir as métricas por classe no TCC: a pergunta é sobre o comportamento do modelo em grupos (as 9 classes), não sobre exames isolados.
\end{citacard}

\begin{citacard}{\pagechip{p. 15} §4 Limitações}
\begin{origquote}“While we identified saliency as the most reliable attribution method through a defined sanity check, comparing it against integrated gradients, LRP, and Grad-CAM, we acknowledge the limitations inherent in relying on a single attribution method for the proposed experiments.”\end{origquote}
\begin{tradquote}“Embora tenhamos identificado a \textit{saliency} como o método de atribuição mais confiável por meio de um teste de sanidade definido, comparando-a com integrated gradients, LRP e Grad-CAM, reconhecemos as limitações de depender de um único método nos experimentos propostos.”\end{tradquote}
\explab{Explicação}
O TCC também usa um só método (GradientSHAP). A ressalva vale igualmente: as conclusões por classe são sobre o par modelo + GradientSHAP, não sobre “a explicação” em geral.
\end{citacard}

\section{Relação com o TCC}
É o trabalho mais próximo do TCC na agregação: atribuições somadas por derivação $\times$ segmento do batimento, que é a mesma ideia das 36 regiões (12 derivações $\times$ pré-QRS / QRS / ST-T), com segmentação mais fina (24 segmentos). Também oferece evidência de que métodos do tipo gradiente $\times$ entrada tendem a concentrar relevância no QRS, o que ajuda a interpretar a concentração do mapa aleatório no QRS observada no eMPRT. Por outro lado, o teste de sanidade do artigo exige treinar modelos de regressão com alvo conhecido, o que o TCC não faz; e o artigo não usa métricas de fidelidade nem compara a qualidade das explicações com o desempenho por classe. A concordância glocal com critérios clínicos é qualitativa e restrita a quatro patologias do PTB-XL.

\section{Como citar no TCC}
\begin{itemize}
  \item Wagner et al.\ (2024) mostraram, com modelos treinados para regredir amplitudes de ondas do ECG, que IG, LRP e Grad-CAM concentram a relevância no complexo QRS mesmo quando a tarefa depende da onda P ou T, o que evidencia um viés dos métodos de atribuição para as regiões de maior amplitude.
  \item A agregação das atribuições por derivação e segmento do batimento segue a proposta de análise “glocal” de Wagner et al.\ (2024).
\end{itemize}

\end{document}
